% Generated draft: requires mathematical and source review.
\documentclass[11pt]{article}
\usepackage[margin=25mm]{geometry}
\usepackage{fontspec}
\setmainfont{Latin Modern Roman}
\usepackage{amsmath,amssymb,mathtools,unicode-math}
\setmathfont{Latin Modern Math}
\usepackage{enumitem,needspace,xurl,hyperref}
\setlength{\parindent}{0pt}
\setlength{\parskip}{5pt}
\setlength{\emergencystretch}{3em}
\newcommand{\N}{\mathbb{N}}
\newcommand{\Z}{\mathbb{Z}}
\newcommand{\Q}{\mathbb{Q}}
\newcommand{\R}{\mathbb{R}}
\newcommand{\C}{\mathbb{C}}
\newcommand{\F}{\mathbb{F}}
\newcommand{\A}{\mathbb{A}}
\newcommand{\PP}{\mathbb{P}}
\newcommand{\E}{\mathbb{E}}
\newcommand{\eps}{\varepsilon}
\newcommand{\dd}{\,\mathrm{d}}
\begin{document}
\section*{UnsolvedMath ID: 30006359}
\textbf{Source problem code:} catalogue entry
\subsection{Definitions and mathematical statement}
Let \(V\) denote the algebraic variety defined by the fusion consistency equations referred to in the supplied source, and let \(\mathcal{D}_{\mathrm{forest}}(V)\) denote its intended forest decomposition. The source record does not supply the equations, the ground field, the ambient parameter space, or the construction of this decomposition, so these symbols designate source-dependent objects rather than fully specified mathematical data.

A nonempty closed subset \(Z\) of a topological space is irreducible if every expression \(Z=Z_1\cup Z_2\), with \(Z_1,Z_2\) closed in \(Z\), has \(Z=Z_1\) or \(Z=Z_2\). An irreducible component is a maximal irreducible closed subset. Write \(\operatorname{Irr}(V)\) for the collection of irreducible components of \(V\), using the intended Zariski topology. The ground field and whether components are to be taken after extension to an algebraic closure remain unspecified.

\textbf{Source conjecture, in schematic form.} Interpreting the forest decomposition as a collection of subsets of \(V\), the intended assertion is
\[
\mathcal{D}_{\mathrm{forest}}(V)=\operatorname{Irr}(V).
\]
Thus each forest piece must be a maximal irreducible closed subset of \(V\), and every irreducible component must occur as a forest piece. This formulation does not determine whether the source actually intends the pieces themselves or their closures; that distinction requires clarification from the missing construction.

The accompanying source qualifications must also be retained. Auxiliary variables whose indices are corrupted in the supplied extraction, printed as \texttt{akl kl ij , bij}, take values in \(\{-1,0,1\}\) determined by \(\Gamma_{\bullet}\). The quantities printed as \texttt{Nijk} are integer linear combinations of variables \(d_{\ell}\), with coefficients determined by \(\Gamma_{\bullet}\). The source states that all consistency equations reduce to linear equations except for the degree-two associativity equation. Neither \(\Gamma_{\bullet}\), the index ranges, the coefficients, nor the associativity equation is defined in the record. Consequently, the displayed conjecture is an explicit statement of the intended target, but a complete instance or quantified family cannot be specified from the supplied evidence alone.
\subsection{Short English statement}
The conjecture says that the pieces obtained from the source's forest construction are exactly the irreducible components of the variety of fusion consistency equations. The record omits the construction, equations, ground field, and allowed input data, so this remains a schematic statement requiring human mathematical review. The supplied record labels the problem open; that status is not independently verified here.
\subsection{Sources}
\begin{itemize}
\item Supplied local registry record: permanent UnsolvedMath ID \(30006359\), title \emph{Irreducible Forest Decomposition of Consistency-Equation Varieties}, source problem code \texttt{catalogue entry}, in the Oberwolfach Reports Open Problems collection. No individual proposer is named; the supplied proposed year is 2025.
\item Supplied report citation: \emph{Subfactors and Applications (2026)}, Oberwolfach Reports, DOI \texttt{10.4171/owr/2025/34}, pp.\ 1823--1886. Supplied source link: \url{https://doi.org/10.4171/owr/2025/34}. The citation's parenthetical year 2026 differs from the 2025 report designation; both are preserved without resolving the discrepancy.
\item Supplied literature reference: \emph{Oberwolfach Report 34/2025}, \url{https://ems.press/content/serial-article-files/52239}.
\item Literature checked at: \texttt{2026-08-22}. This is supplied dated curation, not an independently verified assessment. It labels the problem open and reports that the forest-decomposition assertion is presented as a conjecture and that no proof or counterexample was located. The background explicitly requests independent verification of time-sensitive status and cited claims.
\item The registry labels the statement \texttt{corrected\_verified}, with the qualification that only clear formula delimiters, notation, grammar, and typography were repaired and that mathematical content and scope were preserved. This supplied editorial label does not establish mathematical correctness or resolve the missing definitions.
\end{itemize}
\end{document}
